\chapter{Residual Hyperbolic Candidates for the Cable Detection}\label{app:cable_residual_families}

This appendix records the homological filtering verdicts for all $87$ candidate braid families against the cable target Alexander polynomial $\Phi_{12}(t) = t^4 - t^2 + 1$. The data here support Proposition \ref{prop:cable_filter_outcome} and clarify the residual gap in Theorem \ref{thm:cable_detection_partial}.

The structure of the analysis is identical to that of Appendix \ref{app:granny_verdict_tables}: for each family we have a determinant $\det(I - A_{\text{red}})$ and a characteristic polynomial $P(t) = \det(tI - A_{\text{red}})$, both of which depend only on the family (not on the target polynomial). Only the comparison to the target changes: instead of $(t^2-t+1)^2$, we test against $t^4 - t^2 + 1$.

\section{Surviving Families}

In contrast to the granny case (Appendix \ref{app:granny_verdict_tables}), six families survive both filters at non-negative integer parameter values for the cable target; two of these six (Track 0 / 7 Briefing and Track 0 / 10 Quail) survive at infinitely many parameter values, while the remaining four survive only at finitely many. We list these explicitly.

\begin{table}[h]
\centering
\caption{Candidate families surviving both homological filters at non-negative integer parameter values for the cable target $t^4 - t^2 + 1$.}
\label{tab:cable_residual_families}
\begin{tabular}{@{}l l l >{\raggedright\arraybackslash}p{0.28\textwidth} l@{}}
\toprule
Track & Family & $\det(I - A_{\text{red}})$ & Surviving non-negative parameters & $P(t)$ at those values \\
\midrule
Candelabra & 16 Kieve & $(j{+}1)k+(j{+}k{-}2)m+3j-6$ & $(j,k,m) = (0,3,4)$ & $t^4 - t^2 + 1$ \\
& & & $(j,k,m) = (0,7,0)$ & $t^4 - t^2 + 1$ \\
& & & $(j,k,m) = (1,2,0)$ & $t^4 - t^2 + 1$ \\
\midrule
Track 0 & 7 Briefing & $-(j{+}3)k+(j{-}k)m+j-1$ & infinitely many integer triples & $t^4 - t^2 + 1$ \\
\midrule
Track 0 & 8 Briefing & $km - 1$ & $(k,m) = (1,2)$ & $t^4 - t^2 + 1$ \\
& & & $(k,m) = (2,1)$ & $t^4 - t^2 + 1$ \\
\midrule
Track 0 & 10 Quail & $jk-(j{+}k{+}2)m-1$ & infinitely many integer triples & $t^4 - t^2 + 1$ \\
\midrule
Track 0 & 13 Treasure & $-(k{+}3)m+2k+2$ & $(k,m) = (2,1)$ & $t^4 - t^2 + 1$ \\
\midrule
Track 8 & 9 FAT & $k - 1$ & $k = 2$ & $t^4 - t^2 + 1$ \\
\bottomrule
\end{tabular}
\end{table}

For family Track 0 \,/\, 7 Briefing, the equation
\[
-(j+3)k + (j-k)m + j - 1 = 1
\]
admits infinitely many non-negative integer solutions; representative values include $(j,k,m) = (1,0,1), (2,0,0), (2,1,5), (3,2,11), \dots$ — at each, $P(t) = t^4 - t^2 + 1$.

For family Track 0 \,/\, 10 Quail, the equation
\[
jk - (j+k+2)m - 1 = 1
\]
likewise admits infinitely many non-negative integer solutions; representative values include $(j,k,m) = (1,2,0), (2,1,0), (2,6,1), (3,12,2), \dots$ — at each, $P(t) = t^4 - t^2 + 1$.

\section{Excluded Families}

The remaining $81$ families (that is, $87 - 6$, the six residual families comprising four with finitely many surviving parameter values and two with infinitely many) are excluded by the same Filter 1 / Filter 1* / Filter 2 verdict as in the granny case. The detailed table is identical in structure to Appendix \ref{app:granny_verdict_tables}, with the verdict re-evaluated against the cable target. The excluded families' polynomial mismatches with $t^4 - t^2 + 1$ may be verified directly from the characteristic polynomials recorded there. In particular, the new families $1.1$ Jehu, $1.2$ Jehu, $1.1'$, and $1.2'$ are excluded by Filter 1 against any target (the determinant is constant or admits no non-negative integer assignment giving $|\det| = 1$), so they do not contribute to the residual list.

\section{Discussion}

The residual list above represents the precise gap in Theorem \ref{thm:cable_detection_partial}: any hypothetical hyperbolic knot $K$ with $\widehat{HFK}(K) \cong \widehat{HFK}(C_{2,1}(T_{2,3}))$ must, by the analysis of Chapter \ref{chap:cable_detection}, have its associated 6-braid lying in one of these six families at one of the surviving parameter values. For a complete cable detection theorem, each of these (family, parameter) pairs must be ruled out by an additional argument.

The candidate strategies for closing this gap are discussed in Section \ref{sec:cable_residual_discussion}.